\documentclass[10pt,a4paper]{amsart}
\usepackage[margin=19mm]{geometry}
\usepackage{amsmath,amssymb,mathtools}
\usepackage[scr=rsfs,cal=euler]{mathalfa}
\usepackage{xfrac}
\usepackage[unicode,hidelinks,bookmarks=false]{hyperref}
\newcommand{\ie}{i.e.\ }
\newcommand{\cf}{cf.\ }
\newcommand{\resp}{respectively\ }
\newcommand{\Subsection}{\subsection}
\providecommand{\nobreakdash}{\nobreak\hbox{-}}
\setlength{\emergencystretch}{2em}
\multlinegap0pt
\usepackage{amssymb}
\usepackage{mathtools}
\usepackage[shortlabels]{enumitem}
\usepackage{tikz}
\usepackage{xcolor}
\newtheorem{assumption}{Assumption}
\newtheorem{theorem}{Theorem}
\newtheorem{definition}{Definition}
\newtheorem{proposition}{Proposition}
\newtheorem{lemma}{Lemma}
\newcommand{\E}{\mathbb{E}}
\newcommand{\F}{\mathscr{F}}
\renewcommand{\P}{\mathbb{P}}
\newcommand{\R}{\mathbb{R}}
\newcommand{\ce}{\coloneqq}
\newcommand{\dd}{\ensuremath{\mathrm{d}}}
\newcommand{\dt}{\,\mathrm{d}t}
\newcommand{\dx}{\,\mathrm{d}x}
\newcommand{\grad}{\nabla}
\newcommand{\loc}{\textup{loc}}
\newcommand{\nH}{\ensuremath{\,_\nu H}}
\newcommand{\rmd}{\mathrm{d}}
\newcommand{\sB}{\mathscr{B}}
\newcommand{\sO}{\mathscr{O}}
\newcommand{\sP}{\mathscr{P}}
\title{The stochastic Cahn--Hilliard equation in critical spaces}
\author{Simon Bau, Gideon Chiusole, Sarah Geiss, Katharina Klioba, Tobias Werner}
\date{Source: arXiv:2608.21014v1 (CC BY 4.0)}
\begin{document}
\maketitle

\noindent\emph{Source and licence.} The stochastic Cahn--Hilliard equation in critical spaces. Version arXiv:2608.21014v1 (CC BY 4.0), distributed by arXiv under the Creative Commons Attribution 4.0 International licence, http://creativecommons.org/licenses/by/4.0/, as recorded in the arXiv licence metadata for this exact version and shown on the abstract page https://arxiv.org/abs/2608.21014v1. Full licence notice: \texttt{licenses/arxiv-2608.21014-CC-BY-4.0.txt}.\par\smallskip

\noindent\textit{Editorial scope.} Counted unit: the article's proposition prop:global\_well\_posedness\_dim\_4-part2 (source0.tex lines 1941--1946). The article's complete original proof of it is delivered with it (source0.tex lines 1948--2012, one \texttt{proof} environment of 65 lines). No mathematics is written, repaired or re-derived: the statement and the proof are the article's own lines, in the article's own notation, and no retained line is substituted or re-flowed.  Retained uncounted as the source-local context the proof needs: lines 784--789; lines 873--909; lines 1135--1147; lines 1162--1166; lines 1188--1191; lines 1807--1816; lines 1847--1861; lines 1849--1860; lines 1872--1874; lines 1892--1895; lines 1898--1901; lines 1903--1905.  Nothing was omitted from any retained span, so the package carries no omission record.  No retained line contains a citation, so the wrapper carries no bibliography and no bibliography entry.  No retained line contains a figure environment, an image-inclusion command or drawing code, so no image is delivered and the package declares no asset dependency.  Editorial formatting only, in addition: an AMS \texttt{amsart} preamble that restates the article's own declarations and macros used by the retained lines, namely the environments \texttt{assumption}, \texttt{theorem}, \texttt{definition}, \texttt{proposition}, \texttt{lemma} on separate counters, together with the standard packages those lines need.  The retained environments print in this document as Assumption 1, Theorem 1, Definition 1, Proposition 1, Lemma 1 in order of appearance, whereas the article prints the same items with the numbering of its own full document.  The complete native source of the article as distributed in the arXiv e-print for this exact version is bundled unchanged as \texttt{sources/208218/source0.tex}.
\begin{equation}\tag{AVP}\label{eq:AVP}
    \begin{cases}
        \dd u(t) + Au(t)\dt  &= F(u(t)) \dt + G(t,u(t))\,\dd W(t), \\
        u(0) &=u_0.
    \end{cases}
\end{equation}

\begin{assumption}\label{assumption_weakened_setting}
  Let $d\in \{3,4 \}$. Suppose that 
  \begin{equation*}
      \psi=(\psi_n)_{n\geq 1}:[0,\infty)\times \Omega\times\sO\times \R^{1+d}\to \ell^2,\quad \psi(t,\omega,x,y,z) \ce \psi^0(t,\omega,x,y)+\psi^1(t,\omega,x) \cdot z,
  \end{equation*}
  where $\cdot$ denotes the usual scalar product in $\R^d$ and the functions
 \begin{align*}
     &\psi^0=(\psi^0_n)_{n\geq 1}:[0,\infty)\times \Omega\times\sO\times \R\to \ell^2,\quad
     \psi^1=(\psi^1_n)_{n\geq 1}:[0,\infty)\times \Omega\times\sO\to \ell^2(\R^d)
 \end{align*}
 are $\sP \otimes \sB(\sO)\otimes \sB(\R)$-measurable and $\sP \otimes \sB(\sO)$-measurable, respectively.
 Moreover, we assume that 
\begin{enumerate}
    \item the mappings  $(x,y)\mapsto \psi^0(\cdot,x,y)$ and $x\mapsto \psi^1(\cdot,x)$  are in $C^1(\sO\times \R;\ell^2)$ and $C^1(\sO;\ell^2(\R^d))$  a.e.\ on $[0,\infty)\times \Omega$, respectively;
    \item  there exists a constant $L\geq0$ such that a.e.\ on
    $[0, \infty)\times\Omega\times\sO$, for all $y,y'\in\R$, we have
    \begin{align}
        \|\nabla_x[
            \psi^0(\cdot,y)-\psi^0(\cdot,y')
        ]\|_{\ell^2(\R^d)}
        &\leq L(1+|y|+|y'|)|y-y'|, \label{eq:psi0-gradx-local-lipschitz}
        \\
                \|\partial_y[
            \psi^0(\cdot,y)-\psi^0(\cdot,y')
        ]\|_{\ell^2}
        &\leq L|y-y'|, \label{eq:psi0-dy-lipschitz}
        \\
        \|\psi^0(\cdot,y)\|_{\ell^2}+ \|\nabla_x\psi^0(\cdot,y)\|_{\ell^2(\R^d)}
        &\leq L(1+|y|), \label{eq:psi0-linear-growth}
        \\
        \|\psi^1(\cdot)\|_{\ell^2(\R^d)}
        +
        \|\nabla_x\psi^1(\cdot)\|_{\ell^2(\R^{d\times d})}
        &\leq L. \label{eq:psi1-uniform-bound}
    \end{align}
\end{enumerate}
\end{assumption}

\begin{theorem}[Local well-posedness in the weakened setting] \label{thm:localWP_weakened}
Suppose that Assumption~\ref{assumption_weakened_setting} holds. Then for all $u_0\in L^0_{\F_0}(\Omega;H^1(\sO))$, there exists a (unique) maximal solution $(u,\sigma)$ to \eqref{eq:AVP} such that a.s.\ $u\in C([0,\sigma);H^1(\sO))\cap L^2_{\loc}([0,\sigma);\nH^3(\sO))$. Moreover, the following blow-up criterion holds
\begin{equation}
    \label{eq:blowUp_bare}
    \P \bigg(\sigma < \infty ,\sup\limits_{t\in [0,\sigma)} \| u(t) \|_{H^1(\sO)} + \int_0^\sigma \|u(t) \|^2_{H^3(\sO)} \dt<\infty  \bigg) =0.
\end{equation}

In the noncritical case $d=3$ the blow-up criterion can be relaxed to 
\begin{equation}
    \label{eq:blowUpSubcritd3}
    \P \bigg(\sigma < \infty ,\sup\limits_{t\in [0,\sigma)} \| u(t) \|_{H^1(\sO)} <\infty  \bigg) =0.
\end{equation}
\end{theorem}

\begin{equation}\tag{$\text{SCH}_2$}
\label{eq:SCHtorus}
    \rmd u = [-\Delta^2u+\Delta(u^3-u)
    ]\dt + \sum_{n \geq 1} \psi_n(u,\grad u)\,\rmd W^n,\quad u(0)=u_0,
\end{equation}

    \begin{definition}\label{def:chemPotential}
    The \emph{chemical potential} $\mu$ associated with the solution $u$ of
\eqref{eq:SCHtorus} is given by $\mu\ce -\Delta u+u^3-u$.
    \end{definition}

\begin{proposition} \label{prop:global_well_posedness_dim_4-part1}
	Let $d=4$ and suppose that the assumptions of Theorem~\ref{thm:localWP_weakened} hold. Moreover, let $u_0\in L^0_{\F_0}(\Omega;H^1(\sO))$ and $T>0$. Let $(u,\sigma)$ be the unique maximal solution to \eqref{eq:SCHtorus} and $\mu$ the chemical potential of Definition~\ref{def:chemPotential}. Then, we have
	\begin{equation}\label{eq:prop-global_well_posedness_dim_4-claim1}
	       \sup_{t\in[0,T]\cap [0,\sigma)}\|u(t)\|_{H^1}+\int_0^{\sigma \land T} \|\grad \mu(t)\|_{L^2}^2 \dt < \infty  \qquad \mathbb{P}\text{-a.s.}
	\end{equation}
	If, in addition, $u_0\in L_{\F_0}^4(\Omega;L^4(\sO))\cap L_{\F_0}^2(\Omega;H^1(\sO))$, then
    \begin{equation}\label{eq:prop-global_well_posedness_dim_4-claim2}
    \E\bigg[\sup_{t\in[0,T]\cap[0,\sigma)}  \|u(t)\|^2_{H^1(\sO)} + \int_0^{\sigma \land T} \|\grad \mu(t)\|_{L^2}^2 \dt \bigg] \le C_{\psi,T,\sO}\big(1+\|u_0\|^4_{L^4(\Omega;L^4(\sO))}+\|u_0\|^2_{L^2(\Omega;H^1(\sO))}\big) < \infty.
    \end{equation}
\end{proposition}

\begin{lemma} \label{lem:globalWP_d4_claims}
	Let $d=4$ and suppose that the assumptions of Theorem~\ref{thm:localWP_weakened} hold. Moreover, let $u_0\in L^0_{\F_0}(\Omega;H^1(\sO))$ and $T>0$. Let $(u,\sigma)$ be the unique maximal solution to \eqref{eq:SCHtorus} and $\mu$ the associated chemical potential from Definition~\ref{def:chemPotential}. Then the following three estimates hold a.s.\ for a.e.\ $0\le s <\sigma \land T$
    \begin{align}
        \|\Delta u(s)\|_{L^2}^2  & \le \|\mu(s)\|_{L^2}^2 + 2\| \grad u(s)\|_{L^2}^2 \label{eq:globalWP_d4_eq1}, \\
         \|u^2(s)\nabla u(s)\|_{L^2}^2 & \lesssim
        \|u(s)\|^3_{H^1}
        \left(\|\Delta u(s)\|_{L^2}^2
        +\|u(s)\|^2_{H^1}  \right)
        \left( \|\grad \Delta u(s)\|_{L^2}
        +\|u(s)\|_{H^1}\right),
        \label{eq:globalWP_d4_eq2} \\
        \|\mu(s)\|_{L^2}^4 & \lesssim \|\nabla\mu(s)\|_{L^2}^2 \left(\|u(s)\|_{H^1}^6 + 1\right) + 1
        + \|u(s)\|_{H^1}^{12}. \label{eq:globalWP_d4_eq3}
    \end{align}
\end{lemma}

    \begin{align}
        \|\Delta u(s)\|_{L^2}^2  & \le \|\mu(s)\|_{L^2}^2 + 2\| \grad u(s)\|_{L^2}^2 , \\
         \|u^2(s)\nabla u(s)\|_{L^2}^2 & \lesssim
        \|u(s)\|^3_{H^1}
        \left(\|\Delta u(s)\|_{L^2}^2
        +\|u(s)\|^2_{H^1}  \right)
        \left( \|\grad \Delta u(s)\|_{L^2}
        +\|u(s)\|_{H^1}\right),
         \\
        \|\mu(s)\|_{L^2}^4 & \lesssim \|\nabla\mu(s)\|_{L^2}^2 \left(\|u(s)\|_{H^1}^6 + 1\right) + 1
        + \|u(s)\|_{H^1}^{12}. 
    \end{align}

\begin{equation}\label{eq:muRearranged}
\Delta u = -\mu + u^3 - u,
\end{equation}

\begin{align} \label{eq:ellReg_H2_seminorm}
    |u|_{H^2} &\lesssim \|\Delta u\|_{L^2}+\|u\|_{L^2},\\
    |u|_{H^3} &\lesssim \|\Delta u\|_{H^1}+\|u\|_{L^2} \lesssim \|\grad \Delta u\|_{L^2}+\|\Delta u\|_{L^2}+\|u\|_{L^2},\nonumber
\end{align}

\begin{equation}
\label{eq:proofDivergenceThm}
    \int_{\sO} \Delta u(s) \dx=\int_{\partial\sO} \partial_{\mathbf{n}}u(s) \,\rmd S=0\text{ for a.e.\ }0 \le s <\sigma \land T \text{ a.s.\ }
\end{equation}

\begin{align}
    |u|_{H^3} &\lesssim \|\grad \Delta u\|_{L^2}+\|\Delta u\|_{L^2}+\|u\|_{L^2} \lesssim \|\grad \Delta u\|_{L^2}+\|u\|_{L^2}.\label{eq:ellReg_H3_seminorm}
\end{align}

\begin{proposition} \label{prop:global_well_posedness_dim_4-part2}
	Let $d=4$ and suppose that the assumptions of Theorem~\ref{thm:localWP_weakened} hold. Moreover, let $u_0\in L^0_{\F_0}(\Omega;H^1(\sO))$ and $T>0$. Let $(u,\sigma)$ be the unique maximal solution to \eqref{eq:SCHtorus}. Then we have
	\begin{equation}\label{eq:thm-global_well_posedness_dim_4-claim1}
	  \int_0^{\sigma\wedge T} \|u\|^2_{H^3} \dt < \infty  \qquad \mathbb{P}\text{-a.s.}
	\end{equation}
\end{proposition}

\begin{proof}
It suffices to show $\int_0^{\sigma\wedge T} \|\grad \Delta u\|^2_{L^2} \dt <\infty$ a.s., which can be seen as follows. By the elliptic regularity estimates \eqref{eq:ellReg_H2_seminorm} as well as \eqref{eq:ellReg_H3_seminorm} and absorbing $\|\Delta u\|_{L^2}$ into $\|\grad\Delta u\|_{L^2}$ via \eqref{eq:proofDivergenceThm} and the Poincaré--Wirtinger inequality as before, we obtain
\begin{align}\label{eq:thm:globalWP_d4_H3norm}
    \|u\|_{H^3}^2 = |u|_{H^3}^2+|u|_{H^2}^2 +\|u\|_{H^1}^2
    \lesssim \|\grad\Delta u\|_{L^2}^2+\|\Delta u\|_{L^2}^2+\|u\|_{H^1}^2 \lesssim \|\grad\Delta u\|_{L^2}^2+\|u\|_{H^1}^2.
\end{align} 
Proposition~\ref{prop:global_well_posedness_dim_4-part1} ensures that $\sup_{t\le \sigma \land T} \|u(t)\|_{H^1}<\infty$ a.s. Hence, the integral in time of the second 
term in \eqref{eq:thm:globalWP_d4_H3norm} 
is almost surely finite, so that it suffices to control the integral of $\|\grad\Delta u\|_{L^2}^2$. 

Let $\sigma_m$ be a localising sequence of $(u,\sigma\wedge T)$. We show that $\int_0^{\sigma\wedge T} \|\grad \Delta u\|^2_{L^2} \dt <\infty$ a.s.\ follows from the three estimates \eqref{eq:globalWP_d4_eq1}, \eqref{eq:globalWP_d4_eq2}, and \eqref{eq:globalWP_d4_eq3} of Lemma~\ref{lem:globalWP_d4_claims}.
We use the chemical potential via \eqref{eq:muRearranged} in the first inequality, \eqref{eq:globalWP_d4_eq2} in the second inequality and Young's inequality in the third inequality. We apply \eqref{eq:globalWP_d4_eq1} in the fourth inequality and
\eqref{eq:globalWP_d4_eq3} in the fifth inequality. This results in
\begin{align*}
    & \int_0^{\sigma_m} \|\nabla\Delta u\|_{L^2}^2\dt 
     \le 3\int_0^{\sigma_m} \|\nabla\mu\|_{L^2}^2
    + 9\|u^2\nabla u\|_{L^2}^2
    + \|\nabla u\|_{L^2}^2\dt  \\
    & \le 
    3\int_0^{\sigma_m}
    \|\nabla\mu\|_{L^2}^2
    +
    C\|u\|^3_{H^1}
    \left(\|\Delta u\|_{L^2}^2
    +\|u\|^2_{H^1}  \right)
    \left( \|\grad \Delta u\|_{L^2} +\|u\|_{H^1}\right)+
    \|\nabla u\|_{L^2}^2\dt\\
    & \le  
    3\int_0^{\sigma_m}
    \|\nabla\mu\|_{L^2}^2 \dt +
    C\int_0^{\sigma_m}(1+\|u\|_{H^1}^{10})\dt
    +
    \frac12
    \int_0^{\sigma_m}
    \|\nabla\Delta u\|_{L^2}^2\dt
    +
    C\int_0^{\sigma_m}
    \|u\|_{H^1}^6
    \|\Delta u\|_{L^2}^4\dt  \\
    & \le
    3\int_0^{\sigma_m}
    \|\nabla\mu\|_{L^2}^2 \dt +
    C\int_0^{\sigma_m}(1+\|u\|_{H^1}^{10})\dt
    +
    \frac12
    \int_0^{\sigma_m}
    \|\nabla\Delta u\|_{L^2}^2\dt
    +
    C\int_0^{\sigma_m}
    \|u\|_{H^1}^6
    \|\mu\|_{L^2}^4\dt \\
    & \le
    3\int_0^{\sigma_m}
    \|\nabla\mu\|_{L^2}^2 \dt +
    C\int_0^{\sigma_m}(1+\|u\|_{H^1}^{18})\dt
    +
    \frac12
    \int_0^{\sigma_m}
    \|\nabla\Delta u\|_{L^2}^2\dt
     + C\Big(1+ \sup_{t\leq \sigma_m} \|u\|^{12}_{H^1} \Big)
    \int_0^{\sigma_m}
    \|\nabla\mu\|_{L^2}^2\dt.
 \end{align*}
Here, the constant $C$ changes from line to line. By Proposition~\ref{prop:global_well_posedness_dim_4-part1}, both $\sup_{t\le \sigma_m} \|u(t)\|_{H^1}<\infty$ a.s.\ and $\int_0^{\sigma_m} \|\grad \mu\|_{L^2}^2\dt<\infty$ a.s. Rearranging the terms and bounding all polynomial $H^1$ terms by the pathwise supremum implies the claim.
\end{proof}

\end{document}
